\scanpage{95}
$\mathrm{a)}\Longrightarrow(\mathrm{b)}+\mathrm{c)})$ bien connu (on peut supposer $L$ très ample, et on applique la théorie du pol. de Hilbert)

\[
\begin{aligned}
\mathrm{b)}&\unsure{\longrightarrow|\Longrightarrow}\mathrm{b}')\Longrightarrow\mathrm{b}'')\quad\text{trivial}\\
\mathrm{c)}&\Longrightarrow\mathrm{c}')\Longrightarrow\mathrm{c}'')\quad\text{trivial}
\end{aligned}
\]

On prouve les autres assertions par récurrence sur \deleted{\unsure{b$'$)}} $\dim X=d$, \unsure{tt|tout|cela|c'} est trivial si $d=0$, et bien connu \deleted{\unsure{si}} si $d=1$. Supposons $d\geq2$

$\mathrm{b}'')\Longrightarrow\unsure{\mathrm{c}|\mathrm{d}}'')$ Grâce au

\uline{Lemme 1} \unsure{Si|Si on} les \deleted{conditions b$''$), $\exists$ \unsure{un} $n_0>0$} \inserted{restrictions de $\underline L$ aux sous préschémas \inserted{intègres}} $Z$ de $X$ de $\dim<d$ sont amples, alors $\exists\,n_0$ tel que
\[
 n\geq n_0\Longrightarrow H^i(X,L^{\otimes n})=0\qquad\text{pour }i\geq2\quad\deleted{\text{\unsure{les supp.}}}
\]
\margin{\unsure{par|pour} des restrictions sur $d=\dim X$ Mais il faut supposer $X$ \unsure{projective|projectiv|projectif}.}
% slanted

cf lettre de Mumford.

\deleted{$\mathrm{d)}\Longrightarrow\mathrm{c)}$ trivial}

\deleted{$\mathrm{d)}\Longrightarrow\mathrm{d}')$ d'après le cas déjà \deleted{\unsure{traité}} \inserted{connu} $d=1$.}

$\unsure{\mathrm{c}|\mathrm{d}}'')\Longrightarrow\mathrm{a)}$ On peut supposer (via : EGA III ... que $X$ est intègre. \deleted{\illegible{1}} \deleted{\unsure{Soit}} \deleted{\unsure{f}} $f$ une section de $L^{\otimes n}$, \deleted{\unsure{A}} $D$ \unsure{l'ens|l'un|l'ens.} de ses zéros, alors $D\ne\varnothing$ car $\underline L$ n'est pas trivial, donc $\dim D=d-1$. On prouve

\uline{Lemme} 2 Sous les conditions du lemme 1, supposons que \inserted{$\exists\,n>0$ tel} $\underline L^{\otimes n}$ ait une section. Alors \deleted{$\unsure{\exists}$} $\exists n_0$ tel que $n\geq n_0\Longrightarrow X\overset{!}{\longrightarrow}\mathbb P H^0(X,L^{\otimes n})$ partout définie et non \unsure{constante|const.|constant}

\margin{$d\geq2$}

cf lettre Mumford ....

\uline{Questions} Dans \unsure{les} énoncés d) et d$'$) peut-on supprimer l'hyp d'existence d'une section $\ne0$ ?? \deleted{\unsure{Oui si} \illegible{2}} \deleted{\illegible{1}} Peut-on supprimer l'hyp. que $X$ projectif ??

% archive sheet number: 94
